\documentclass{article}

\usepackage{amsmath,amssymb}
\usepackage{xurl}
\usepackage[hidelinks]{hyperref}

% Shared commands for notes in docs/paper-gaps/.

\DeclareUnicodeCharacter{2080}{\ifmmode{}_{0}\else\textsubscript{0}\fi}
\DeclareUnicodeCharacter{2081}{\ifmmode{}_{1}\else\textsubscript{1}\fi}
\DeclareUnicodeCharacter{2082}{\ifmmode{}_{2}\else\textsubscript{2}\fi}
\DeclareUnicodeCharacter{2099}{\ifmmode{}_{n}\else\textsubscript{n}\fi}

\newcommand{\Lean}{\textsc{Lean}}
\ExplSyntaxOn
\NewDocumentCommand{\leanid}{m}{
  \ifmmode
    \textsf{\detokenize{#1}}
  \else
    \group_begin:
    \urlstyle{sf}
    \tl_set:Nn \l_tmpa_tl {#1}
    \tl_replace_all:Nnn \l_tmpa_tl {\_} {_}
    \exp_args:NV \path \l_tmpa_tl
    \group_end:
  \fi
}
\ExplSyntaxOff
\providecommand{\tr}{\operatorname{tr}}
\newcommand{\tnleanIssueBase}{https://github.com/LionSR/TNLean/issues/}
\newcommand{\tnleanPrBase}{https://github.com/LionSR/TNLean/pull/}
\newcommand{\tnleanDocBase}{https://sirui-lu.com/TNLean/docs/}
\newcommand{\ghissue}[1]{\href{\tnleanIssueBase#1}{GitHub issue \##1}}
\newcommand{\ghpr}[1]{\href{\tnleanPrBase#1}{PR \##1}}
\newcommand{\doclink}[1]{\href{\tnleanDocBase#1.html}{\path{#1}}}

% Dirac notation and the identity operator, as in blueprint/src/macros/common.tex.
% \providecommand keeps note-local definitions authoritative.
\usepackage{dsfont}
\providecommand{\ket}[1]{|#1\rangle}
\providecommand{\bra}[1]{\langle #1|}
\providecommand{\braket}[2]{\langle #1 | #2 \rangle}
\providecommand{\ketbra}[2]{|#1\rangle\!\langle #2|}
\providecommand{\Id}{\mathds{1}}
\providecommand{\one}{\mathds{1}}
\providecommand{\C}{\mathbb{C}}
\providecommand{\MN}[1]{M_{#1}(\C)}
\providecommand{\MD}{M_D(\C)}

% External referencing for paper sources.  The build helper writes sanitized
% auxiliary files under build/paper-aux/ and excludes duplicated source labels.
\usepackage{xr}

% Import a generated paper auxiliary file with a caller-supplied label prefix.
% The candidate paths support builds from docs/paper-gaps/, the repository root,
% and build/paper-aux/ itself.
\newcommand{\importpaperaux}[2]{%
  \IfFileExists{../../build/paper-aux/#2.aux}{%
    \externaldocument[#1]{../../build/paper-aux/#2}%
  }{%
    \IfFileExists{build/paper-aux/#2.aux}{%
      \externaldocument[#1]{build/paper-aux/#2}%
    }{%
      \IfFileExists{#2.aux}{%
        \externaldocument[#1]{#2}%
      }{}%
    }%
  }%
}

% General external reference macros
\newcommand{\paperref}[2]{\ref{#1-#2}}
\newcommand{\papereqref}[2]{(\ref{#1-#2})}

% CPSV16 source (1606.00608 / MPDO-22-12-17-2.tex) external document and shortcuts
\importpaperaux{cpsv16-}{1606.00608/MPDO-22-12-17-2-tnlean}
\newcommand{\cpsvref}[1]{\paperref{cpsv16}{#1}}
\newcommand{\cpsveqref}[1]{\papereqref{cpsv16}{#1}}

% Verdict marker: \gapnote{<kind>}{<status>}. Kind is the policy
% classification (clarification, local-correction, scope-restriction,
% unfaithful, false-source, open-gap); status is open, wip (elimination
% underway), resolved, or historical. Machine-read by the site index;
% prints a verdict line.
\newcommand{\gapnote}[2]{\par\noindent\textbf{Verdict:} #1 (#2).\par}


\title{The Positive Primitive in the Generalized-Cocycle Corollary}
\author{}
\date{2026-08-29}

\begin{document}
\maketitle
\gapnote{local-correction}{resolved}

\section{Source argument}

Let $G$ be finite and let the positive generalized $2$-cocycle $\Omega$ take
values in the embedded subgroup $\mathbb R_{>0}\subset\mathbb C^\times$.
The proof of
\cite[Corollary~\texttt{cor:positive\_trivial}, lines~5950--5960]
{FrancoRubio2025MPUGauging} invokes the preceding identity
\begin{align}
    \Omega^{|G|} &= d\delta,
    \qquad
    \delta^x_g=\det X^x_g,
    \label{eq:fbc25_positive_generalized_power}
\end{align}
and then takes a positive $|G|$-th root.

The determinant cochain $\delta$ need not itself be positive. Indeed, the
regular-action matrix factors as a permutation matrix times a diagonal matrix,
and the permutation determinant can equal $-1$. Thus a positive root of
$\delta^x_g$ is not defined in general.

\section{Positive primitive}

Take the pointwise complex norm before taking the root:
\begin{align}
    \chi^x_{0,g}
        &=\lvert\delta^x_g\rvert^{1/|G|}
          =\lvert\det X^x_g\rvert^{1/|G|}.
    \label{eq:fbc25_positive_generalized_primitive}
\end{align}
Since $\delta^x_g\ne0$, every value of $\chi_0$ lies in
$\mathbb R_{>0}\subset\mathbb C^\times$. Multiplicativity of the complex norm
and positivity of $\Omega$ turn
Eq.~\ref{eq:fbc25_positive_generalized_power} into
\begin{align}
    \Omega^{|G|} &= d\lvert\delta\rvert.
    \label{eq:fbc25_positive_generalized_norm_power}
\end{align}
Positive real powers are multiplicative, so
$(d\chi_0)^{|G|}=\Omega^{|G|}$. Injectivity of the $|G|$-th power on
$\mathbb R_{>0}$ gives $d\chi_0=\Omega$.

\section{Formal treatment}

The formal construction takes the pointwise norm of the determinant cochain
and then its positive $|G|$-th root. The block label $x$ is retained in
Eq.~\ref{eq:fbc25_positive_generalized_primitive}. This changes no hypothesis
or conclusion of the corollary.\footnote{Formal declarations:
\leanid{TNLean.Algebra.LSymbol.positivePrimitive} and
\leanid{TNLean.Algebra.LSymbol.exists\_positive\_coboundary} in
\doclink{TNLean/Algebra/PositiveGeneralizedCocycle}.}

\textbf{Verdict: resolved local correction.}
The norm removes the possible permutation sign while preserving the
coboundary identity required by the positive-root argument.

\bibliographystyle{alpha}
\bibliography{references}

\end{document}